%============================================================
%  TIỂU LUẬN template (Vietnamese essay / thesis-style report)
%  Start one with:  scripts/new_note.sh <Course> essay-vi "Tên môn học" TieuLuan
%  Build:  tectonic <file>.tex     (XeTeX only: Vietnamese comes from fontspec, no babel)
%
%  Choices that matter:
%    - 13pt body: the standard "report" class silently ignores 13pt and falls back
%      to 10pt, so this uses KOMA's scrreprt with fontsize=13pt.
%    - Margins 30/20/25/25 mm, 1.5 line spacing, 1 cm first-line indent: the usual
%      Vietnamese university format. Check your faculty's rules and adjust.
%    - Vietnamese is not hyphenated like English, so hyphenation is switched off.
%    - Put logo.png next to this file (the title page shows a box until you do).
%============================================================
\documentclass[a4paper,fontsize=13pt,oneside,numbers=noenddot,chapterprefix=true]{scrreprt}
\usepackage{fontspec}
\setmainfont{texgyretermes}[Extension=.otf, UprightFont=*-regular, BoldFont=*-bold,
  ItalicFont=*-italic, BoldItalicFont=*-bolditalic]
\usepackage{amsmath,amssymb}
\usepackage[a4paper, left=30mm, right=20mm, top=25mm, bottom=25mm]{geometry}
\usepackage{setspace}
\onehalfspacing
\usepackage{indentfirst}
\setlength{\parindent}{1cm}
\usepackage[bottom]{footmisc}
\usepackage{graphicx}
\usepackage{booktabs}
\usepackage{enumitem}
\usepackage[hidelinks,bookmarks=false]{hyperref}
\usepackage{tocbasic}

% Vietnamese is not hyphenated the way English is.
\hyphenpenalty=10000 \exhyphenpenalty=10000 \tolerance=3000 \emergencystretch=3em

% ---- Vietnamese labels (set at \begin{document}, after KOMA sets its English ones) ----
\AtBeginDocument{%
  \renewcommand{\chaptername}{Chương}%
  \renewcommand{\contentsname}{MỤC LỤC}%
  \renewcommand{\bibname}{TÀI LIỆU THAM KHẢO}%
  \renewcommand{\figurename}{Hình}%
  \renewcommand{\tablename}{Bảng}}
% Headings in the body font (KOMA defaults to sans), chapters centred:
%   CHƯƠNG 1
%   Title of the chapter
\setkomafont{disposition}{\normalfont\bfseries}
\renewcommand*{\chapterformat}{\MakeUppercase{\chaptername}~\thechapter\autodot}
\RedeclareSectionCommand[beforeskip=0pt, afterskip=18pt, innerskip=6pt]{chapter}
\renewcommand*{\raggedchapter}{\centering}
\addtokomafont{chapter}{\Large}
\addtokomafont{chapterprefix}{\large}
\addtokomafont{section}{\large}
\addtokomafont{subsection}{\itshape}
\setcounter{secnumdepth}{3}
\setcounter{tocdepth}{2}

% ---- who and what (edit these) ----
\newcommand{\StudentName}{Your Name}
\newcommand{\StudentID}{YOUR-ID}
\newcommand{\CourseName}{Tên môn học}
\newcommand{\CourseCode}{MÃ MÔN}
\newcommand{\Major}{Ngành học}
\newcommand{\EssayTitle}{TÊN ĐỀ TÀI TIỂU LUẬN}
\newcommand{\University}{TRƯỜNG ĐẠI HỌC \dots}
\newcommand{\Faculty}{KHOA \dots}

%============================================================
\begin{document}

\begin{titlepage}
  \centering
  {\large\bfseries ĐẠI HỌC QUỐC GIA THÀNH PHỐ HỒ CHÍ MINH\par}
  {\large\bfseries \University\par}
  {\large\bfseries \Faculty\par}
  \vspace{1.2cm}
  \IfFileExists{logo.png}{\includegraphics[width=3cm]{logo.png}}{\fbox{\parbox[c][3cm][c]{3cm}{\centering logo.png}}}\par
  \vspace{1.2cm}
  {\LARGE\bfseries TIỂU LUẬN\par}
  \vspace{4pt}
  {\large Môn học: \CourseName\ (\CourseCode)\par}
  \vspace{1cm}
  {\Large\bfseries \EssayTitle\par}
  \vfill
  \begin{tabular}{@{}ll@{}}
    Sinh viên thực hiện: & \StudentName\\
    Mã số sinh viên: & \StudentID\\
    Ngành: & \Major\\
  \end{tabular}
  \vfill
  {\large Thành phố Hồ Chí Minh, năm \the\year\par}
\end{titlepage}

\pagenumbering{roman}
\tableofcontents
\clearpage
\pagenumbering{arabic}

\chapter*{MỞ ĐẦU}
\addcontentsline{toc}{chapter}{MỞ ĐẦU}
\section*{1. Lý do chọn đề tài}
\addcontentsline{toc}{section}{1. Lý do chọn đề tài}
Nội dung \dots
\section*{2. Mục đích và nhiệm vụ nghiên cứu}
\addcontentsline{toc}{section}{2. Mục đích và nhiệm vụ nghiên cứu}
\section*{3. Đối tượng và phạm vi nghiên cứu}
\addcontentsline{toc}{section}{3. Đối tượng và phạm vi nghiên cứu}
\section*{4. Phương pháp nghiên cứu}
\addcontentsline{toc}{section}{4. Phương pháp nghiên cứu}
\section*{5. Kết cấu của tiểu luận}
\addcontentsline{toc}{section}{5. Kết cấu của tiểu luận}

\chapter{Cơ sở lý luận}
\section{Khái niệm}
Trích dẫn theo kiểu tác giả (năm) \cite[tr.~187]{mac}.

\chapter{Thực trạng và vận dụng}
\section{Nội dung}

\chapter*{KẾT LUẬN}
\addcontentsline{toc}{chapter}{KẾT LUẬN}
Tóm tắt kết quả \dots

\begin{thebibliography}{99}
\addcontentsline{toc}{chapter}{TÀI LIỆU THAM KHẢO}
\item[] \textbf{Tài liệu tiếng Việt}
\bibitem{mac} Tác giả (2021), \emph{Tên sách}, Nxb. Chính trị quốc gia, Hà Nội.
\item[] \textbf{Tài liệu tiếng Anh}
\bibitem{en1} Author, A. (2020), \emph{Book Title}, Publisher, City.
\end{thebibliography}

\end{document}
